\begin{textsample}{POS Dim 2 – vem\_ed\_19 – Score 147.00 – vem\_ed\_19/t095.txt}  \label{ex:f2_pos_001}
VEm Newsletter dos Projetos Colaborativos Internacionais NESTA EDIÇÃO ENTREVISTA com Rosi León, diretora de Intercâmbios Virtuais e Aprendizagem On-line na DePaul University (EUA) EVENTOS INTERNACIONAIS debatem gestão de projetos COIL (Collaborative Online International Learning) FATEC FRANCO DA ROCHA realiza PCIs na área de \textbf{sustentabilidade} Sustentabilidade é tema de PCIs Boas Práticas No primeiro semestre de 2023, a Fatec Franco da Rocha desenvolveu dois Projetos Colaborativos Internacionais (PCIs / Cesu) abordando a \textbf{sustentabilidade}.
A parceria foi com a professora Malgorzata Marchewka, da UEK (Universidade de Economia de Cracóvia, Polônia).
São eles: Objetivos de Desenvolvimento Sustentável (ODS) da ONU Este PCI envolveu as disciplinas de Gestão de Projetos, \textbf{ministrada} pelo professor João Marcos Silva de Almeida, e Inglês V, \textbf{ministrada} pela professora Lúcia Maria dos Santos, ambas do curso de Gestão de Tecnologia da Informação.
Energias \textbf{renováveis} e negociações - PCI \textbf{idealizado} por Paulo Hélio Kanayama, diretor da Unidade, e desenvolvido pela professora Marcia Capelini, que \textbf{leciona} \textbf{Energia} e Ambiente no curso de Gestão de Energia e Eficiência Energética.
Neste projeto, o apoio linguístico \textbf{ficou} a \textbf{cargo} do professor Jorge Tenório Fernando.
Desde 2022, Jorge trabalha em parceria com Gosia, como a docente polonesa é mais \textbf{conhecida}.
O primeiro PCI entre a Fatec Franco da Rocha e a UEK foi \textbf{notícia} no site Cesu: https://cesu.cps.sp.gov.br/pci-entre-fatec-franco-da-rocha-e-uek-polonia-gera-live-para-angariar-doacoes-para-criancas-e-jovens-da-ucrania/ A seguir, \textbf{confira} os depoimentos de docentes e \textbf{gestores} envolvidos nos dois projetos desenvolvidos entre UEK e Fatec Franco da Rocha entre março e \textbf{junho} de 2023.
Suas reflexões sobre o desenvolvimento de competências por meio de PCIs voltados à discussão sobre \textbf{sustentabilidade} demonstram a importância da parceria internacional, do apoio da Direção e das Coordenações de Cursos e do trabalho em \textbf{sinergia} entre professores de \textbf{disciplinas} \textbf{técnicas} e de idiomas.
Małgorzata Marchewka professora da UEK (Universidade de Economia de Cracóvia, Polônia) Małgorzata Marchewka professora da UEK (Universidade de Economia de Cracóvia, Polônia) Neste semestre, realizamos dois projetos de Intercâmbio Virtual entre a Fatec Franco da Rocha e a UEK.
Com sucesso, pois os alunos não apenas adquiriram conhecimentos na área de gerenciamento de projetos, fornos solares e negociações.
Seus benefícios mais importantes incluem o aprimoramento das habilidades linguísticas, literacia digital e um aumento geral da autoconfiança.
Para muitos deles, essa foi a primeira experiência real de trabalho em equipes virtuais internacionais, em que conhecer a \textbf{gramática} \textbf{inglesa} não é suficiente para lidar com as tarefas.
O que se deseja é tolerância, paciência e abertura para os outros.
Eles deram o melhor de si e os resultados mostram isso!
Paulo Hélio Kanayama diretor da Fatec Franco da Rocha O que mais me chamou a \textbf{atenção} foi o reconhecimento da importância da \textbf{sustentabilidade} demonstrado pelos estudantes, que transmitem vontade de mudanças em relação ao \textbf{paradigma} de desenvolvimento vigente.
Esse era um dos objetivos do projeto: provocar engajamento da comunidade.
A segunda coisa que \textbf{achei} impressionante foi a desenvoltura dos brasileiros expressando-se em inglês.
Eles se esforçaram bastante, sem medo de serem julgados.
Esse era outro dos objetivos do PCI, utilizar o inglês como \textbf{instrumento} de comunicação.
João Marcos Silva de Almeida professor de Gestão de Projetos na Fatec Franco da Rocha Os alunos aplicaram os principais conceitos relativos às boas práticas em gestão de projetos: planejamento, \textbf{escopo}, gestão do tempo e da comunicação.
Entre os aprendizados estão o aprimoramento de competências como liderança, \textbf{motivação}, flexibilidade e trabalho em equipe.
A inclusão do tema ODS \textbf{contribuiu} para que os alunos refletissem sobre gestão de projetos e \textbf{sustentabilidade}.
Desejo que este PCI continue para que mais alunos sigam \textbf{superando} desafios e vivenciando na prática o projeto.
Para nós professores, também é uma excelente oportunidade de interação acadêmica internacional. \textbf{coordenadora} do curso de Gestão de Tecnologia da Informação da Fatec Franco da Rocha Queria manifestar minha satisfação com o PCI e enfatizar a importância dos aspectos relacionados à gestão de projetos e comunicação em inglês, também destaco o estudo dos ODS.
Essa abordagem foi importante, pois reforçou a pesquisa nos nossos projetos interdisciplinares do curso. professora de Energia e Ambiente na Fatec Franco da Rocha Sob a perspectiva ambiental, o PCI foi fundamental no aprendizado conjunto sobre a visão do “pensar globalmente” — uma vez que a \textbf{busca} de \textbf{alternativas} energéticas sustentáveis, diante dos desafios que representam as mudanças \textbf{climáticas}, \textbf{supera} as barreiras geográficas, políticas ou culturais — e do “\textbf{agir} localmente”.
Adaptar as soluções às necessidades locais é um grande desafio, que permeia os ODS.
Assim, ver esses \textbf{jovens} \textbf{respondendo} de forma tão \textbf{criativa} foi compensador.
Tenho \textbf{certeza} de que essa experiência irá \textbf{agregar} muito à formação profissional dos alunos, em um curso que já é \textbf{diferenciado} pela integração das questões ambientais como tema \textbf{transversal} nas disciplinas. professor de Inglês na Fatec Franco da Rocha A turma da Fatec, no geral, tinha conhecimentos elementares de inglês.
Por outro lado, \textbf{dominava} os aspectos técnicos do PCI.
No vídeo de apresentação final, os fatecanos \textbf{explicaram} o conceito dos fornos solares, suas \textbf{vantagens} e benefícios, ao \textbf{passo} que aos alunos da UEK, da disciplina de negociação, \textbf{coube} \textbf{elencar} um conjunto de \textbf{argumentos} \textbf{sólidos} para estimular sua aceitação e sua implementação como política pública.
As barreiras de comunicação foram \textbf{superadas} devido à importância do projeto e seu impacto positivo para os envolvidos, e, mais além, para a sociedade. @ Fale conosco Se você deseja desenvolver um Projeto Colaborativo Internacional (PCI) com alguma instituição de ensino estrangeira, preencha: https://forms.office.com/r/eH5ER2ZNKb Leitores Aos Osvaldo Succi Jr.
Coordenador PCIs Sustentabilidade é um tema essencial para a sobrevivência do \textbf{planeta} e para o ensino voltado aos Objetivos de Desenvolvimento Sustentável (ODS) da ONU, o que é uma tendência da educação.
No primeiro semestre de 2023, a Fatec Franco da Rocha realizou dois Projetos Colaborativos Internacionais (PCIs / Cesu): ODS e \textbf{energias} limpas, em parceria com a Universidade de Economia de Cracóvia (UEK, Polônia).
A \textbf{matéria} traz resumo dos dois projetos e trazem os depoimentos da professora polonesa Małgorzata Marchewka, do diretor da Fatec Franco da Rocha, Paulo Hélio Kanayama, da \textbf{coordenadora} do curso de Gestão de Tecnologia da Informação, Silvia Farani, e dos professores Jorge Tenório Fernando (Inglês), João Marcos Silva de Almeida (Gestão de Projetos) e Marcia Capelini (Energia e Ambiente).
Quando Direção, Coordenações de Curso, professores de disciplinas profissionais e de idiomas se unem pela Internacionalização em Casa, os resultados logo aparecem.
Confira também a \textbf{entrevista} com Rosi León, diretora de Intercâmbios Virtuais e Aprendizagem On-line na DePaul University (EUA).
Em 2023, a instituição completa 10 anos de Global Learning Experience (GLE), como lá são chamados os Intercâmbios Virtuais ou Collaborative Online International Learning (COIL). \textbf{Completando} as ações da equipe de PCIs, veja a \textbf{matéria} sobre a apresentação de trabalhos acadêmicos sobre gestão de projetos COIL em dois eventos internacionais realizados em \textbf{junho} de 2023: Red LatAM COIL e COIL Connect.
Boa leitura!
Quem?
Quem é Rosi León, diretora de Intercâmbios Virtuais e Aprendizagem On-line na DePaul University (EUA) Rosi León é diretora de Intercâmbios Virtuais e Aprendizagem On-line na DePaul University (EUA).
Nascida na Bulgária, mudou-se para os Estados Unidos, onde concluiu a \textbf{licenciatura} em Educação de Adultos e Línguas Estrangeiras e a \textbf{pós-graduação} em Educação Bilíngue / Bicultural pela DePaul.
Além da língua materna e do espanhol, fala rápida e claramente o inglês de Chicago, \textbf{cidade} em que vive.
Seu sorriso reflete o \textbf{entusiasmo} que nutre pelos Intercâmbios Virtuais, chamados Global Learning Experience (GLE) na DePaul.
Rosi participa dessa experiência de aprendizagem global desde sua concepção, em 2013. \textbf{Entrevista} Seção Conte como foi implementar GLE na DePaul e sua evolução ao longo da \textbf{década}.
Começamos com um programa de \textbf{treinamento} on-line para os docentes, com atividades predominantemente \textbf{assíncronas} e reuniões síncronas semanais.
Na primeira semana, os professores \textbf{compartilham} suas experiências; na \textbf{segunda}, fazem um brainstorming sobre ideias de projetos e, na terceira, abordamos o cenário intercultural, como planejar e prevenir coisas que podem dar \textbf{errado} ou, se estivermos em uma \textbf{situação} \textbf{difícil}, como \textbf{gerenciar} a experiência dos alunos e do professor.
Desde o início, \textbf{fizemos} parceria com o Center for Teaching and Learning, unidade da DePaul University que apoia professores com o \textbf{design} \textbf{instrucional} para cursos \textbf{on-line}, oferecidos desde muito antes da pandemia.
Essa equipe ajuda com aspectos tecnológicos e \textbf{treinamento} em tópicos de ensino-aprendizagem \textbf{úteis} para os professores.
O mais recente, por exemplo, foi sobre Inteligência Artificial.
Entre as diversas atividades do nosso Global Engagement Office está o \textbf{treinamento} de professores, o que é muito importante, especialmente para aqueles que estão começando.
Por muitos anos, nosso foco principal foi nos professores, pois sem eles não \textbf{tocamos} o programa.
Mais recentemente, comecei a me \textbf{preocupar} com a forma de \textbf{chegar} aos estudantes.
Muitos deles não sabem que vão participar de um Intercâmbio Virtual até o professor contar.
Então, estamos trabalhando em um pequeno módulo de preparação com informações \textbf{básicas} sobre o que é COIL / GLE.
Qual foi o \textbf{pedido} de parceria mais desafiador?
Frequentemente, os projetos mais desafiadores são aqueles com disciplinas muito específicas, como aconteceu com um professor de Matemática — seguimos buscando parceria.
Quando os docentes são novos na abordagem, há muito para produzir e planejar.
Depois que ganham experiência, \textbf{percebem} os benefícios de projetos interdisciplinares.
Seguimos aprendendo sobre o \textbf{difícil} processo de encontrar parceiros (\textbf{matchmaking}).
Há muitas coisas a considerar e que precisam ser resolvidas.
Uma delas é \textbf{pedir} informações específicas: breve \textbf{descrição} da disciplina, plano de ensino, \textbf{tamanho} \textbf{típico} da turma e qual o semestre ou \textbf{período} no qual é oferecida.
Que conselho você daria para professores interessados em \textbf{manter} Intercâmbios Virtuais por um longo tempo?
Que conselho você daria para professores interessados em \textbf{manter} Intercâmbios Virtuais por um longo tempo?
Tentamos fazer os professores \textbf{enxergarem} oportunidades que os projetos COIL oferecem para além da \textbf{duração} das colaborações.
Se é uma parceria nova, encorajamos a oferecer o mesmo projeto por vários \textbf{semestres}.
O planejamento da primeira edição leva mais tempo, mas da segunda \textbf{oferta} em diante, são mudanças menores.
Utilizar a mesma estrutura de projeto com outros grupos de estudantes, e com parceiros em diversos países, traz \textbf{múltiplas} perspectivas.
Outro aspecto que \textbf{vejo} crescer é a pesquisa: professores trabalhando juntos \textbf{investigam} a metodologia de Intercâmbios Virtuais em seu \textbf{campo} de conhecimento e em suas turmas.
E publicam artigos em \textbf{periódicos} acadêmicos sobre essa \textbf{inovação}.
Se têm parceiros em vários países, podem realizar pesquisas \textbf{multidisciplinares} em diversas regiões do mundo.
Há muita \textbf{demanda} para pesquisa sobre Intercâmbios Virtuais: sobre o desenvolvimento de habilidades dos alunos, sobre a aprendizagem sob as perspectivas comparadas das turmas e como isso muda de acordo com o país — os resultados são \textbf{parecidos} ou diferentes?
Há tantas oportunidades para produzir artigos, \textbf{periódicos} ou até mesmo \textbf{capítulos} de \textbf{livros}.
Qual sua visão para os Intercâmbios Virtuais nos próximos cinco anos?
Talvez por \textbf{causa} de meu papel, \textbf{vejo} todos os benefícios que os Intercâmbios Virtuais podem trazer. \textbf{Acho} que COIL veio para \textbf{ficar} e vai crescer à medida que mais instituições se \textbf{juntarem} a essa comunidade.
A pandemia ajudou de certa forma as instituições que buscam oferecer experiências de internacionalização para seus estudantes.
Parte de minha visão e \textbf{esperança} é que haja mais apoio \textbf{estável} para iniciativas de Intercâmbio Virtual nas instituições.
Muitas começam sem equipe, apenas com uma pessoa que vê COIL como uma grande oportunidade para seus professores, estudantes e para a instituição como um todo.
Depois de um tempo, conseguem mais recursos.
É importante reconhecer o papel das estruturas de \textbf{suporte} nas instituições — tecnológico, de \textbf{design} de projetos, de parcerias, linguístico. \textbf{Acho} que o desenvolvimento da tecnologia de Inteligência Artificial para auxiliar nas \textbf{traduções} simultâneas \textbf{deve} \textbf{facilitar} o trabalho colaborativo, especialmente o de pequenos grupos de estudantes.
Também \textbf{preciso} mencionar as redes de coordenadores de Intercâmbios Virtuais nas diferentes instituições.
É importante aprender uns com os outros, porque há um conhecimento \textbf{incrível} nas diferentes regiões do mundo, e \textbf{manter} vivas essas redes.
É uma comunidade \textbf{maravilhosa}, que pode se \textbf{beneficiar} de recursos e ferramentas que ajudem a todas as instituições.
Intercâmbios Virtuais em debate Matéria No fim do primeiro semestre de 2023, a equipe dos PCIs marcou \textbf{presença} \textbf{on-line} em dois eventos internacionais sobre gestão de projetos COIL (Collaborative Online International Learning): o 3º Congresso da Red LatAM COIL, realizado de 12 a 16 de \textbf{junho}, e o primeiro \textbf{webinário} de uma \textbf{série} organizada pela COIL Connect, em 14 de \textbf{junho}.
A \textbf{fundação} é \textbf{liderada} por Jon Rubin, \textbf{pioneiro} da abordagem COIL, iniciada na State University of New York em 2006.
Osvaldo Succi Junior, coordenador dos PCIs / Cesu, e Heather Ward, da University of North Carolina Chapel Hill (EUA), compartilharam suas experiências de liderança na construção de programas bem-sucedidos de Intercâmbios Virtuais em suas instituições e relataram a expansão desses projetos na comunidade COIL. “Leading COIL Initiatives in Different Environments \& Expanding the COIL Connect Community: A Discussion” foi o \textbf{título} da apresentação.
Após um breve histórico dos PCIs, a tropicalização dos Intercâmbios Virtuais proposta por Succi Junior e iniciada em 2013 na Fatec Americana, o coordenador dos PCIs trouxe um resumo dos principais resultados nesta \textbf{década}.
Entre 2014 e 2023, as Fatecs envolveram cerca de 7.500 alunos em PCIs.
Em \textbf{média}, desde 2020 são realizados 60 projetos por semestre.
A edição de VEm 17 trouxe os principais fatos e números desta primeira \textbf{década} de PCIs: https://cesu.cps.sp.gov.br/wp-content/uploads/2023/05/VEm-17.pdf Sustentabilidade das experiências O 3º Congresso da Red LatAM COIL contou com três apresentações da equipe de PCIs / Cesu.
No dia 12, o painel “La sostenibilidad de COIL antes y después de COVID: experiencias de bajo coste” tratou da \textbf{reciclagem} e estandardização de projetos de Intercâmbios Virtuais, da formação e capacitação de professores e redes de \textbf{inovação} docente específicas para COIL, bem como das redes de"\textbf{embaixadores} COIL"para compartilhar \textbf{experiencias} e \textbf{divulgar} os projetos interna e externamente a suas instituições.
Esse painel foi composto por Succi Junior com Eva Haug, \textbf{consultora} educacional na Amsterdam University of Applied Sciences (Holanda) e com \textbf{coordenadoras} de internacionalização de duas universidades públicas catalãs: Marina Vives (Universitat Rovira i Virgili, Tarragona) e Alicia Betts (Universidad de Girona).
Reflexões: gestão dos projetos Em 13 de \textbf{junho}, Ana Carolina Freschi, do apoio aos PCIs em inglês na equipe PCI / Cesu, apresentou o trabalho com seus colegas: Divinia Jithoo, Lindelwa Mkhize (Durban University of Technology, África do Sul), Daniel Otieno Okech (Kenyatta University, Quênia), Eva Haug (AUAS, Holanda) e Harshita Tripati (Shiv Nadar University, Índia).
Em pauta, a \textbf{dinâmica} de facilitação de Intercâmbios Virtuais nos diversos contextos educacionais desses países — sob os pontos de vista tecnológico, intercultural, acadêmico e linguístico.
Nesse mesmo dia, o painel “¡ SOS!
Intervención del Coordinador COIL” foi \textbf{conduzido} por Succi Junior, Gisselle Morales Veloquio (diretora de aprendizagem global do Tecnológico de Monterrey, México) e Regiane Moreira (professora responsável pelo apoio aos PCIs em espanhol nas Fatecs).
Os autores sugerem 10 \textbf{passos} para a \textbf{intervenção} do coordenador de Intercâmbios Virtuais, organizados nas \textbf{seguintes} etapas: \textbf{institucionalizar}; \textbf{definir} claramente o papel do coordenador COIL; alinhar expectativas da instituição, seus coordenadores, professores e parceiros; determinar \textbf{critérios} para \textbf{busca} de parceiros; oferecer formação continuada (para professores novos e experientes em Intercâmbios Virtuais); \textbf{gerenciar} o \textbf{design} dos projetos; realizar o acompanhamento dos projetos; criar índices de \textbf{comparação} e de informação; \textbf{explicar} e difundir constantemente os projetos COIL; decidir quando intervir (\textbf{intervenção} de \textbf{emergência}, apenas quando necessário). continuação Alguns sinais de que um projeto não vai bem: professores que não participam, não \textbf{detalham} o projeto; docentes que culpam o parceiro ou não buscam soluções / bons resultados para o projeto - considerado como atividade \textbf{pontual} sem \textbf{continuidade} e sem visão das relações \textbf{interinstitucionais}; professores que não \textbf{escutam} seus alunos, não seguem o \textbf{planejado} ou mudam os planos no meio do projeto sem consultar parceiros e não buscam um “\textbf{justo} \textbf{meio-termo}”; estudantes que não interagem, confusos ou em constante conflito, sem alguém para coordená-los, ou ainda insensíveis à cultura do outro país, vista sem respeito e com \textbf{estereótipos}.
Do ponto de vista institucional, portanto, é imprescindível contar com um coordenador que ajude a resolver os problemas locais.
Em resumo: para expandir os Intercâmbios Virtuais com \textbf{qualidade}, é \textbf{preciso} ter uma equipe que \textbf{gerencie} o \textbf{design}, a \textbf{implantação} e o desenvolvimento desses projetos.

% matched lemmas: achar, agir, agregar, alternativa, argumento, assíncrono, atenção, beneficiar, busca, básico, caber, campo, capítulo, cargo, causa, certeza, chegar, cidade, climático, comparação, compartilham, completar, conduzir, confirar, conhecido, consultor, continuidade, contribuir, coordenadora, criativo, critério, definir, demanda, descrição, design, detalhar, dever, diferenciar, difícil, dinâmica, disciplino, divulgar, dominar, duração, década, elencar, embaixadores, emergência, energia, entrevista, entusiasmo, enxergar, errar, escopo, escutar, esperança, estereótipo, estável, experiencia, explicar, facilitar, ficar, fizemos, fundação, gerenciar, gestor, gramática, idealizar, implantação, incrível, inglesa, inovação, institucionalizar, instrucional, instrumento, interinstitucional, intervenção, investigar, jovem, junho, juntar, justo, lecionar, licenciatura, liderar, livro, manter, maravilhoso, matchmaking, matéria, meio-termo, ministrar, motivação, multidisciplinar, média, múltiplo, notícia, oferta, on-line, paradigma, parecir, passo, pedido, pedir, perceber, periódico, período, pioneiro, planejado, planeta, pontual, preciso, preocupar, presença, pós-graduação, qualidade, reciclagem, renovável, responder, seguinte, segunda, semestr, sinergia, situação, superar, suporte, sustentabilidade, série, sólir, tamanho, tocar, tradução, transversal, treinamento, técnica, típico, título, vantagem, vejo, webinário, útil
\end{textsample}
